\documentclass[11pt]{article}

\usepackage[a4paper,margin=1in]{geometry}
\usepackage{amsmath,amssymb}
\usepackage{graphicx}
\usepackage{booktabs}
\usepackage{xcolor}
\usepackage{verbatim}
\usepackage[colorlinks=true,linkcolor=blue!55!black,citecolor=blue!55!black,
            urlcolor=blue!55!black]{hyperref}

% smaller monospace for all code blocks
\makeatletter
\renewcommand{\verbatim@font}{\normalfont\ttfamily\small}
\makeatother

% worked-example delimiters (robust: rules + bold lead-in, no fragile boxes)
\definecolor{democ}{RGB}{40,90,120}
\newcommand{\demohead}[1]{\par\smallskip\noindent%
  \textcolor{democ}{\rule{\linewidth}{0.7pt}}\par\nobreak\vspace{2pt}%
  \noindent\textbf{\textcolor{democ}{Worked example.} #1}\par\nobreak\vspace{2pt}}
\newcommand{\demofoot}{\par\vspace{2pt}\noindent%
  \textcolor{democ}{\rule{\linewidth}{0.7pt}}\par\smallskip}

\newcommand{\R}{\mathbb{R}}
\newcommand{\E}{\mathbb{E}}
\newcommand{\code}[1]{\texttt{#1}}

\title{\vspace{-1.2cm}\bfseries How the Recursive Global Solver Works\\[4pt]
\large Chebyshev interpolation, Smolyak grids, and the two-branch expectation
in \code{code/global/solver\_recursive/}}
\author{Computational methods note}
\date{\today}

\begin{document}
\maketitle

\begin{abstract}
\noindent This note explains, from the ground up and in plain terms, the
three numerical ideas your recursive global solver in
\code{code/global/solver\_recursive/} is built from. \textbf{Chebyshev
interpolation} lets us replace an unknown policy function by a short list of
numbers. \textbf{Smolyak sparse grids} let us do that in six state
dimensions without the point count exploding. And the \textbf{two-branch
expectation} in \code{point\_map.py} is how sovereign-default risk actually
enters the banker's optimality conditions: the future is split into a
no-default and a default continuation, weighted by a priced probability. I
keep the maths self-contained and lead with intuition, with small worked
examples you can check by hand, and a dedicated section shows exactly which
functions get approximated and where the interpolation plugs into the solve
loop. The last two sections are the ones you asked
for specifically: how the ``branch'' is computed in your code (and why an
earlier representative-branch approximation, now removed, got the sign
wrong), and a careful analysis of \emph{why the solve fails at the
$\mu{=}2$ grid while $\mu{=}1$ works.}
\end{abstract}

\tableofcontents

\vspace{0.4cm}
\section{Orientation: what runs where}
\label{sec:orient}

The \code{code/global/} tree now has a \textbf{single} solving machinery, and
this note is about it.

\begin{itemize}
\item \code{solver\_recursive/} --- the \textbf{recursive global} solver, and
      the only one. Rather than a single time path it solves for \emph{decision
      rules}: functions that map any state $S$ to the equilibrium choices,
      represented with Chebyshev polynomials on a Smolyak grid, with expectations
      done by genuine multi-branch quadrature. \code{main.py} is the driver over
      it: it solves the steady state once and runs the TFP shock, the sovereign-risk
      pass-through, and the OMT/TPI activation comparison, each a full time-iteration
      solve. \textbf{This is the code this note explains.}
\end{itemize}

An earlier perfect-foresight solver (\code{solver\_pf/}, since removed) stacked
the whole transition into one big Newton system and priced sovereign risk with a
single \emph{representative} post-default economy; that approximation produced an
\emph{expansionary} response to sovereign risk --- the wrong sign. The recursive
solver, using the real distribution over the default event, recovers the right
sign (risk lowers output and consumption). We will see exactly why in
\S\ref{sec:branch}.


\section{What we are actually solving for}
\label{sec:what}

Start with the object we want. In a recursive equilibrium the ``answer'' is
not a single number or a single time path --- it is a set of \emph{rules},
one for each equilibrium quantity, each a function of the current state.
Your state is the six-vector defined in
\code{state\_grid.py} (\code{STATE\_NAMES}):
\[
S \;=\; (\,K_D,\;K_F,\;P_D,\;P_F,\;B_D,\;s\,),
\]
the two capital stocks, the two banks' gross deposit obligations, the stock
of risky $D$-government debt, and the risk factor $s$ that sets the priced
default probability. For each state the solver must return values for the
seven ``market-clearing'' unknowns (\code{decision\_rules.py},
\code{SOLVE7}),
\[
(N_D,\,N_F,\,K'_D,\,K'_F,\,r^{\text{dep}}_D,\,r^{\text{dep}}_F,\,p),
\]
plus a handful of read-off objects (banker valuations $\alpha$, bond prices
$Q_b$, consumption $C$). Each of these is a \emph{function of the whole
state} --- a surface over a six-dimensional box.

Why can we not just linearize (Taylor-expand) around the steady state, as a
perturbation method would? Two reasons. First, the banks face a leverage
constraint that binds only \emph{sometimes} --- it is slack in normal times
and clamps down in a crisis. That switch puts a \emph{kink} in the true
policy, and a Taylor polynomial around one point cannot see a kink somewhere
else. Second, the entire point of the exercise is the behaviour \emph{away}
from the steady state, in the stressed region. So we need a \emph{global}
method: one that represents the whole surface, not just its slope at a
point.

The recipe for that is the \emph{projection method}: pick a flexible family
of functions (here, Chebyshev polynomials), and choose the member of that
family that makes the model's equations hold exactly at a finite set of
sample states (the \emph{collocation nodes}). Everything below is machinery
for doing that well. Three questions drive it:
\emph{which functions and which sample points} (\S\ref{sec:cheb}),
\emph{how to place sample points in six dimensions without needing millions
of them} (\S\ref{sec:smolyak}), and
\emph{how to compute the expectations the equations contain}
(\S\ref{sec:branch}).


\section{Approximating a function of one variable}
\label{sec:cheb}

Forget economics for a moment. Suppose you have a smooth function $f$ on the
interval $[-1,1]$ and you want to store it in a computer using only a few
numbers, so that you can later evaluate it anywhere, cheaply and accurately.
The plan: write $f$ as a weighted sum of fixed ``building-block''
polynomials,
\[
f(x)\;\approx\; p_n(x)\;=\;c_0\,T_0(x)+c_1\,T_1(x)+\dots+c_n\,T_n(x),
\]
and store just the weights $c_0,\dots,c_n$. The building blocks $T_k$ are the
\emph{Chebyshev polynomials}, and the reason they are the right choice is the
subject of this section.

\subsection{The Chebyshev polynomials}

Here is the slickest definition. On $[-1,1]$ every $x$ can be written as
$x=\cos\theta$ for some angle $\theta\in[0,\pi]$. Define
\[
T_k(x)\;=\;\cos(k\theta)\qquad\text{where }x=\cos\theta.
\]
That looks like trigonometry, not a polynomial --- but it \emph{is} a
polynomial of degree $k$ in $x$, because $\cos(k\theta)$ can always be
rewritten using powers of $\cos\theta = x$. You do not have to take this on
faith; the following recurrence builds them and is exactly the code in
\code{state\_grid.py}:
\begin{equation}
T_0(x)=1,\qquad T_1(x)=x,\qquad
T_{k}(x)=2x\,T_{k-1}(x)-T_{k-2}(x).
\label{eq:recur}
\end{equation}

\demohead{Building the first few Chebyshev polynomials by hand.}
Start from $T_0=1$ and $T_1=x$ and turn the crank of \eqref{eq:recur}:
\[
\begin{aligned}
T_2 &= 2x\cdot x - 1 = 2x^2-1,\\
T_3 &= 2x(2x^2-1) - x = 4x^3-3x,\\
T_4 &= 2x(4x^3-3x) - (2x^2-1) = 8x^4-8x^2+1.
\end{aligned}
\]
Sanity check against the trig definition: $\cos(2\theta)=2\cos^2\theta-1$, and
substituting $x=\cos\theta$ gives $T_2=2x^2-1$. They agree. Notice each
$T_k$ stays between $-1$ and $+1$ on the interval (it is a cosine), and it
wiggles more as $k$ grows --- $T_k$ has exactly $k$ roots in $[-1,1]$. Those
wiggles are what let a sum of them trace out any shape.
\demofoot

\noindent In code, \eqref{eq:recur} fills a table of $T_0,\dots,T_n$ at a
list of points:
\begin{verbatim}
def chebyshev_basis_1d(x, max_deg):
    T = np.empty((x.size, max_deg + 1))
    T[:, 0] = 1.0
    if max_deg >= 1:
        T[:, 1] = x
    for k in range(2, max_deg + 1):
        T[:, k] = 2.0 * x * T[:, k - 1] - T[:, k - 2]
    return T
\end{verbatim}
Why use the recurrence instead of just storing $f$ as $a_0+a_1x+a_2x^2+\dots$
(ordinary powers)? Because the powers $1,x,x^2,x^3,\dots$ start to look
almost identical to each other on $[-1,1]$ as the degree grows (they all
flatten out), so solving for the coefficients $a_k$ becomes numerically
unstable --- tiny rounding errors blow up. The Chebyshev polynomials, being
disguised cosines, stay ``spread out'' and keep the computation stable.

\subsection{Why the sample points matter: the Runge disaster}
\label{sec:runge}

To pin down $n+1$ coefficients we need the fit to match $f$ at $n+1$ sample
points (nodes). The obvious choice --- evenly spaced points --- is
surprisingly terrible. The textbook illustration uses
$f(x)=1/(1+25x^2)$, a perfectly smooth, innocent-looking bump. Fit a
polynomial through evenly spaced points and, as you add more points, the fit
does not get better --- it develops wild oscillations near the ends of the
interval that grow \emph{without bound} (Figure~\ref{fig:runge}, left). This
is the \emph{Runge phenomenon}.

\begin{figure}[h]
\centering
\includegraphics[width=\textwidth]{runge.pdf}
\caption{The same degree-14 fit of the same smooth bump. On evenly spaced
nodes (left) the polynomial swings violently near $\pm1$; on Chebyshev nodes
(right) it tracks the function. Only the placement of the dots differs.}
\label{fig:runge}
\end{figure}

The cure is to \emph{cluster} the sample points toward the ends of the
interval. The right clustering is given by the \emph{Chebyshev nodes} (used
throughout \code{state\_grid.py}):
\begin{equation}
x_k \;=\; -\cos\!\Big(\frac{\pi k}{m-1}\Big),\qquad k=0,\dots,m-1.
\label{eq:extrema}
\end{equation}
There is a lovely picture behind this formula (Figure~\ref{fig:nodes}):
space $m$ points evenly \emph{around} a semicircle, then drop each straight
down onto the horizontal axis. The shadows they cast are the Chebyshev
nodes, and because a semicircle is steep near its ends, the shadows bunch up
near $\pm1$ --- exactly the clustering that defeats Runge.

\begin{figure}[h]
\centering
\includegraphics[width=0.72\textwidth]{nodes.pdf}
\caption{The Chebyshev nodes are evenly spaced angles on a semicircle,
projected onto the line. The projection crowds them toward the endpoints.}
\label{fig:nodes}
\end{figure}

There is a precise sense in which this choice is nearly the best possible. A
number called the \emph{Lebesgue constant} measures how much worse
interpolation can be than the theoretically ideal fit. For evenly spaced
nodes this number grows \emph{exponentially} with the degree (catastrophe);
for Chebyshev nodes it grows only like $\log n$ --- so slowly that in
practice interpolating on Chebyshev nodes is about as good as it is possible
to do. You do not need the formula; the takeaway is: \emph{Chebyshev nodes
turn a potentially explosive problem into a tame one.}

\subsection{Fitting is just solving a small linear system}
\label{sec:fit}

Say we have decided on the degrees $0,\dots,n$ and the nodes
$x_0,\dots,x_n$. ``Fit $f$'' means: find the coefficients $c$ so that
$p_n(x_k)=f(x_k)$ at every node. Writing that out for each node gives $n+1$
linear equations in the $n+1$ unknowns $c$,
\begin{equation}
\underbrace{\begin{pmatrix}
T_0(x_0)&\cdots&T_n(x_0)\\
\vdots& &\vdots\\
T_0(x_n)&\cdots&T_n(x_n)
\end{pmatrix}}_{\text{basis matrix }\Phi}
\begin{pmatrix}c_0\\ \vdots\\ c_n\end{pmatrix}
=
\begin{pmatrix}f(x_0)\\ \vdots\\ f(x_n)\end{pmatrix},
\qquad\text{i.e. } \Phi\,c=f.
\label{eq:collocation}
\end{equation}
The matrix $\Phi$ depends only on the grid, never on $f$. So the code
computes it once, factorizes it once (an ``LU factorization,'' just
bookkeeping that makes repeated solves fast), and reuses it forever after.
In \code{state\_grid.py} that factor is \code{self.\_lu}; the two workhorse
methods are \code{fit} (given values $f$, solve for $c$) and \code{eval}
(given $c$, evaluate $p_n$ anywhere). Every rule in
\code{decision\_rules.py} is stored this way, as a vector of point values
plus its fitted coefficients.

\demohead{A three-point fit you can do on paper.}
Take three Chebyshev nodes on $[-1,1]$: from \eqref{eq:extrema} with $m=3$
they are $x=-1,\,0,\,+1$. The building blocks are $T_0=1$, $T_1=x$,
$T_2=2x^2-1$. Evaluate them at the three nodes to get the basis matrix, and
match a function with values $a,b,c$ at $x=-1,0,1$:
\[
\Phi=\begin{pmatrix}1&-1&1\\ 1&0&-1\\ 1&1&1\end{pmatrix},
\qquad \Phi\begin{pmatrix}c_0\\ c_1\\ c_2\end{pmatrix}
=\begin{pmatrix}a\\ b\\ c\end{pmatrix}.
\]
Solving the $3\times3$ system by hand gives the tidy answer
$c_0=\tfrac{b}{2}+\tfrac{a+c}{4}$, $c_1=\tfrac{c-a}{2}$,
$c_2=\tfrac{a+c}{4}-\tfrac{b}{2}$.
Try it on $f(x)=e^{x}$, whose values are $a=e^{-1}=0.368$, $b=1$,
$c=e=2.718$. You get $p_2(x)=1.000+1.175\,x+0.543\,x^2$. This reproduces
$e^x$ exactly at the three nodes and to within a few percent in between
(worst error about $0.07$ near $x=\tfrac12$). That is the whole idea in
miniature: \emph{three numbers stand in for a function}. Add a couple more
Chebyshev degrees and the error drops to machine precision --- which is the
next point.
\demofoot

\subsection{How accurate is it? Smoothness decides}
\label{sec:accuracy}

The accuracy of a Chebyshev fit depends almost entirely on how \emph{smooth}
the target function is. Two regimes matter (Figure~\ref{fig:conv}):

\begin{itemize}
\item If $f$ is \textbf{very smooth} (infinitely differentiable, no kinks),
      the error falls \emph{geometrically} --- each extra degree multiplies
      it by a constant factor well below one. This is called \emph{spectral
      accuracy}: you reach machine precision with a couple of dozen degrees.
      The smooth curve in Figure~\ref{fig:conv} is a straight line on a log
      scale, the signature of geometric decay.
\item If $f$ has a \textbf{kink} (a corner, where the slope jumps), the
      error falls only \emph{algebraically} --- painfully slowly, like
      $1/n$ or $1/n^2$ --- and the fit develops small ripples near the
      corner (the same family of wiggles as the Runge oscillations, here
      called Gibbs oscillations).
\end{itemize}

\begin{figure}[h]
\centering
\includegraphics[width=0.62\textwidth]{convergence.pdf}
\caption{Largest fit error versus polynomial degree, on Chebyshev nodes. A
smooth target (blue) races down to machine precision; a kinked target (red)
crawls. The occasionally-binding leverage constraint puts a kink in the
model's true policy, so the red curve is the one to keep in mind.}
\label{fig:conv}
\end{figure}

This is not a minor footnote for us: the bank leverage constraint is
occasionally binding, which puts a genuine kink into the true policy
functions. So in the region where the constraint switches on, a polynomial
basis is fighting its worst case. That single fact is behind much of
\S\ref{sec:risks} and all of \S\ref{sec:mu2}.

\subsection{The one thing you must never do: extrapolate}
\label{sec:extrap}

A Chebyshev fit is trustworthy \emph{inside} $[-1,1]$. Outside it, the
polynomials do not just get a bit worse --- they explode. A degree-$n$
Chebyshev polynomial grows roughly like $(\,|x|+\sqrt{x^2-1}\,)^{n}$ once
$|x|>1$, which is enormous.

\demohead{How fast extrapolation blows up.}
The degree-10 block $T_{10}$ satisfies $T_{10}(1)=1$ at the edge of the
interval. Step just outside:
\[
T_{10}(1.05)\approx 12,\qquad
T_{10}(1.2)\approx 252,\qquad
T_{10}(1.5)\approx 7{,}560.
\]
A $20\%$ overshoot of the box edge magnifies that one building block by a
factor of $250$; if its coefficient is not tiny, the ``fit'' returns
nonsense --- a negative consumption, say, which then feeds into an
expectation and produces \code{NaN}. (These numbers are computed directly
from the recurrence \eqref{eq:recur}.)
\demofoot

\noindent This is why the recursive solver \emph{clips} every evaluation
point back into the box before interpolating. You see it in
\code{point\_map.py} wherever the continuation is evaluated, e.g.
\code{cont.eval\_all(d\_next, cont.grid.clip(Sn))} --- \code{clip} projects
any stray next-period state back onto the box boundary, so the rule is never
asked to extrapolate. It is a deliberate, visible guard, not a silent
patch.


\section{Six dimensions without the explosion: Smolyak grids}
\label{sec:smolyak}

Section \ref{sec:cheb} handled one variable. Our state has six. The naive
extension is a \emph{tensor grid}: put $m$ Chebyshev nodes on each axis and
take every combination. That is $m^6$ points. Even a coarse $m=5$ gives
$5^6=15{,}625$ collocation points --- and at every one of them the solver
must solve a nonlinear system and compute an expectation, on every iteration.
Push to $m=9$ and it is over half a million. This geometric blow-up with the
number of dimensions is the famous \emph{curse of dimensionality}
(Figure~\ref{fig:growth}, grey line). Tensor grids are hopeless past three
or four states.

\begin{figure}[h]
\centering
\includegraphics[width=0.6\textwidth]{growth.pdf}
\caption{Points needed versus number of states. The full tensor grid
explodes geometrically; the Smolyak grids grow gently. At our six states the
Smolyak grid needs $13$ points ($\mu{=}1$), $85$ ($\mu{=}2$) or $389$
($\mu{=}3$) --- against tens of thousands for a comparable tensor grid.}
\label{fig:growth}
\end{figure}

\subsection{The idea: keep the cheap combinations, drop the expensive ones}

Smolyak's insight is that most of a tensor grid is wasted. The points that
resolve fine detail in \emph{many} dimensions \emph{at once} (the deep
interior of the grid) are hugely numerous but contribute little; the useful
points resolve detail in one or two dimensions at a time. Smolyak keeps only
the latter. Concretely it works with \emph{levels}: level-$i$ resolution in
one dimension uses $m_i$ nested Chebyshev nodes,
\[
m_1=1,\quad m_i=2^{\,i-1}+1\ \ (i\ge2)\ \Longrightarrow\ m=1,3,5,9,\dots
\]
``Nested'' means every node of one level is reused at the next --- level 2's
three points contain level 1's single point, and so on (\code{\_level\_points}
in \code{state\_grid.py}). A combination of per-dimension levels is written
as a multi-index $\mathbf{i}=(i_1,\dots,i_d)$. A full tensor grid would allow
every combination; Smolyak keeps only the \emph{cheap} ones, those whose
total level is small:
\begin{equation}
\text{keep }\mathbf{i}\quad\Longleftrightarrow\quad
\sum_{j=1}^{d}(i_j-1)\;\le\;\mu.
\label{eq:smolyak}
\end{equation}
The single number $\mu$ (``mu'' in the code) is the \emph{approximation
level}: $\mu=0$ is one point, $\mu=1$ adds one level of detail along each
axis, $\mu=2$ adds pairwise interactions, and so on. This is exactly the
\code{\_multi\_indices} routine, which enumerates the kept combinations.

\demohead{Building the $\mu=2$ grid in two dimensions --- all 13 points.}
Take $d=2$, $\mu=2$. Rule \eqref{eq:smolyak} keeps every $\mathbf{i}=(i_1,i_2)$
with $(i_1-1)+(i_2-1)\le 2$. There are six such combinations. For each, we
take the tensor product of the \emph{new} points that level introduces
(level 1 adds $\{0\}$; level 2 adds the endpoints $\{-1,1\}$; level 3 adds
$\{\pm0.707\}$):
\[
\begin{array}{lll}
(1,1)\to\{0\}\times\{0\} & 1\text{ point (the centre)} &\\
(1,2),(2,1)\to\{0\}\times\{-1,1\}\ \&\ \{-1,1\}\times\{0\} & 4\text{ points (axis ends)} &\\
(1,3),(3,1)\to\{0\}\times\{\pm.707\}\ \&\ \{\pm.707\}\times\{0\} & 4\text{ points (axis interior)} &\\
(2,2)\to\{-1,1\}\times\{-1,1\} & 4\text{ points (the corners)} &
\end{array}
\]
Total: $1+4+4+4=13$ points --- which is exactly what
\code{SmolyakGrid([-1,-1],[1,1],mu=2).n} returns. Figure~\ref{fig:blocks}
plots them, coloured by which combination they came from. Notice the shape:
a dense cross along the axes plus the four corners, with the expensive dense
\emph{interior} of a full grid simply absent. That gap is the saving.
\demofoot

\begin{figure}[h]
\centering
\includegraphics[width=0.56\textwidth]{smolyak_blocks.pdf}
\caption{The 13 points of the two-dimensional $\mu=2$ Smolyak grid, coloured
by the multi-index block that contributes them (the worked example above).
This is the actual \code{state\_grid.SmolyakGrid} construction.}
\label{fig:blocks}
\end{figure}

The same rule in higher dimensions and higher levels gives the sparse
crosses of Figure~\ref{fig:2d}. The bookkeeping detail that makes it all
work: points and polynomial degrees are added in lock-step
(\code{\_new\_points} alongside \code{\_new\_degrees}), so the basis matrix
$\Phi$ is square and invertible, and ``fit'' is again a single solve of
$\Phi c=f$ --- just with 13, 85, or 389 rows instead of 3.

\begin{figure}[h]
\centering
\includegraphics[width=\textwidth]{smolyak2d.pdf}
\caption{Full tensor grid (left) versus Smolyak at $\mu=2$ and $\mu=3$. The
sparse grids keep the axes and a few cross terms and discard the dense
middle.}
\label{fig:2d}
\end{figure}

\subsection{What $\mu$ buys, and what it costs}

A higher $\mu$ represents more interactions between states: $\mu=1$ is
essentially \emph{separable} (it can bend along each axis but captures almost
no cross-effects --- how the response along one state depends on another);
$\mu=2$ adds all pairwise cross-terms; $\mu=3$ adds triples. More faithful,
but more points (13, 85, 389 at $d=6$), and --- as \S\ref{sec:mu2} is
entirely about --- more of those points sit in awkward corners of the state
space. Two more features of \code{state\_grid.py} matter later:

\begin{itemize}
\item \code{build\_state\_box} places the box around the steady state, with
      per-state half-widths (\code{k\_band}, \code{p\_band}, \code{b\_band})
      and the $s$ range set so the priced default probability spans about
      $0.01\%$ to $3\%$. The bands are deliberately narrow for $K$ (capital
      barely moves) and wider for $P,B$ (the movers).
\item \code{mu\_vec} allows an \emph{anisotropic} grid --- a different level
      per state. You can ask for cross-terms only in the dimensions where
      the constraint boundary actually moves ($s,P,B$) and keep the nearly
      fixed capital dimensions at level 1. This is the natural lever to try
      against the $\mu=2$ failure, and \S\ref{sec:mu2} comes back to it.
\end{itemize}


\section{Connecting the tools to the model: what we fit, and how well}
\label{sec:connect}

Sections \ref{sec:cheb}--\ref{sec:smolyak} were deliberately abstract: how to
approximate ``a function $f$,'' and how to place sample points in many
dimensions. This section ties that back to the model --- \emph{which}
functions we actually fit, \emph{why} the fit is needed at all, and
\emph{how well} it converges for the particular functions at hand. It is the
bridge between the machinery and the economics, and it sets up the
$\mu=2$ analysis in \S\ref{sec:mu2}.

\subsection{The functions we fit are the decision rules}

The things we run Chebyshev over are the model's \emph{decision rules} ---
the equilibrium quantities written as functions of the state. The code
(\code{decision\_rules.py}) lists them in two groups:
\begin{itemize}
\item \code{SOLVE7} --- the seven market-clearing unknowns
      $N_D,N_F,K'_D,K'_F,r^{\text{dep}}_D,r^{\text{dep}}_F,p$;
\item \code{DERIVED} --- eight quantities read off the optimality conditions:
      the banker valuations $\alpha_D,\alpha_F$, the bond prices
      $Q_{bD},Q_{bF}$, consumption $C_D,C_F$, and deposit holdings
      $A_D,A_F$.
\end{itemize}
That is $15$ quantities. Each is stored \emph{per default regime}
$d\in\{0,1\}$ (\code{RuleSet}), so we are really approximating
$15\times2=30$ separate scalar functions, each one a surface over the
six-dimensional state box. Every one is represented the same way: its values
at the Smolyak points, plus the Chebyshev coefficients fitted to them.

Here is a concrete way to feel the compression. At $\mu=2$ each rule is
$85$ numbers, so the \textbf{entire} global equilibrium of the model is just
$30\times85=2{,}550$ numbers (at $\mu=1$ it is $30\times13=390$). Everything
downstream --- every impulse response, every stressed-state allocation --- is
reconstructed from those numbers by evaluating the interpolant. The solver's
whole job is to find the right $2{,}550$ numbers.

\subsection{The problem we solve is a fixed point in those rules}

What makes finding them non-trivial is that the rules are self-referential.
The equilibrium conditions at a state $S$ involve \emph{next period's} rules,
because a bank choosing today must forecast what it will do tomorrow. So the
rules we are solving for appear on both sides: the rules agents use to look
ahead must equal the rules that describe their behaviour. That is a
\emph{fixed-point} problem, and \code{recursive\_main.time\_iteration} solves
it by repeated substitution (Figure~\ref{fig:loop}): guess the rules, treat
them as tomorrow's behaviour (the ``continuation''), re-solve every grid
point today, refit, and repeat until nothing moves.

\begin{figure}[h]
\centering
\includegraphics[width=0.86\textwidth]{loop.pdf}
\caption{The time-iteration loop, with the two places the Chebyshev/Smolyak
machinery is load-bearing marked in colour: evaluating the frozen rules at
off-grid next-states $S'$ inside the expectation (\emph{red}), and refitting
solved point values into coefficients (\emph{green}).}
\label{fig:loop}
\end{figure}

\subsection{Why the interpolant is needed at all}

This is the heart of the connection, and the single best answer to ``how do
Chebyshev and Smolyak plug into the model.'' When we solve the equilibrium at
one grid-point state $S$ (Figure~\ref{fig:loop}, right box), the banker's
expectation looks one step ahead to the states $S'$ that today's choices lead
to --- one $S'$ for each Gauss--Hermite shock draw, in each default branch
(\S\ref{sec:branch}). Those next-states are \textbf{almost never grid points};
they are arbitrary interior states. To ask ``what will banks do at $S'$?'' we
must read the continuation rules \emph{off the grid} --- and that is exactly,
and only, what a global interpolant provides. In the code it is the one line
\code{cont.eval\_all(d', cont.grid.clip(Sn))}: a matrix product of the
Chebyshev basis at $S'$ against the stored coefficients, with \code{clip}
guarding the extrapolation cliff of \S\ref{sec:extrap}.

So the three ingredients earn their places precisely here. Without a
\emph{functional} representation you would know the rules only at the $85$
grid points and could not form the expectation at the reachable $S'$ at all.
\emph{Chebyshev} makes reading between the grid points accurate and cheap.
\emph{Smolyak} keeps the number of grid points affordable in six dimensions,
so that ``solve at every point, every iteration'' is feasible. The two
coloured callouts in Figure~\ref{fig:loop} are the two spots this matters:
interpolating the continuation off-grid, and the refit.

\subsection{One fit serves all thirty rules}

A practical efficiency worth noting: because the basis matrix $\Phi$ depends
only on the grid, the \emph{same} factorized $\Phi$ (the LU factor of
\S\ref{sec:fit}) fits every rule. \code{grid.fit} takes a whole array of
point values and returns coefficients in a single solve, and
\code{eval\_all} evaluates all rules at a point with one basis matrix. The
per-iteration cost is therefore dominated by the $85\times2$ pointwise
nonlinear solves, not by the fitting --- the interpolation layer is nearly
free once $\Phi$ is factorized.

\subsection{How well it converges --- for these particular functions}
\label{sec:connect-conv}

Finally, ``convergence'' means two different things here, and keeping them
apart is essential for \S\ref{sec:mu2}.

\begin{enumerate}
\item \textbf{Approximation convergence} --- how close the \emph{best}
      Chebyshev fit of a given rule gets to the true rule, as we raise the
      level $\mu$. This is the smoothness story of \S\ref{sec:accuracy}
      applied to the decision rules. Over most of the state box --- where the
      leverage constraint is comfortably slack or comfortably binding --- the
      rules are smooth, and the fit is near-spectral: a modest $\mu$ already
      captures them and raising it drives the error down fast. But every one
      of these rules has a \emph{kink} where the constraint switches on (the
      $\max\{\cdot,0\}$ in the closed-form multiplier). Right at that crease
      the rule is only continuous, not smooth, so --- exactly like the red
      curve in Figure~\ref{fig:conv} --- the approximation converges slowly
      there and a higher-degree fit rings. In one line: \emph{the functions
      we fit are mostly smooth with one crease, and the crease is where all
      the trouble concentrates.}
\item \textbf{Solver convergence} --- whether the time-iteration loop
      actually \emph{reaches} its fixed point. This is a different question.
      Even if each rule \emph{could} be represented well, the loop might not
      get there: it might oscillate, or settle on the wrong equilibrium. The
      $\mu=1$ grid converges in this second sense; the $\mu=2$ grid does not.
\end{enumerate}

\noindent The honest consequence, flagged earlier: at $\mu=1$ the magnitudes
are only \emph{indicative}, because a near-separable fit cannot capture how
the responses interact across states --- but the \emph{sign and mechanism}
are trustworthy, because every rule it uses is read at well-behaved,
converged points. Sharpening the magnitudes needs a fit that resolves both
the crease and the cross-state curvature. That is exactly what $\mu=2$ was
supposed to buy --- and \S\ref{sec:mu2} is the story of why, here, it cannot.


\section{The numerical risks, in plain terms}
\label{sec:risks}

Before the branch and the $\mu=2$ analysis, here is the short list of ways
this kind of solver goes wrong, and the guard for each. They all reappear in
\S\ref{sec:mu2}.

\begin{enumerate}
\item \textbf{Extrapolation (\S\ref{sec:extrap}).} Evaluating a rule outside
      its box returns garbage. Guard: \code{grid.clip} before every
      interpolation in \code{point\_map.py}.
\item \textbf{Kinks and Gibbs wiggles (\S\ref{sec:accuracy}).} The
      occasionally-binding constraint is a kink; a global polynomial fits it
      only slowly and rings near it. In the recursive solver the multiplier
      is handled by Bocola's \emph{closed form},
      $\mu=\max\{1-\E[\Omega]R\,n/(\lambda\cdot\text{assets}),\,0\}$
      (\code{point\_map.py}), capped just below 1 so it stays finite when a
      deep-crisis net worth goes small. (The perfect-foresight solver in
      \code{solver\_pf/} instead uses Fischer--Burmeister smoothing; same
      economics, different numerical dressing.)
\item \textbf{Bad corners poisoning a global fit.} Because the Chebyshev
      basis is \emph{global}, every coefficient depends on \emph{all} node
      values (that single $\Phi c=f$ solve). So if the solve fails at a few
      awkward corner states, freezing them at stale values still corrupts
      the fitted surface \emph{everywhere}. This is the central villain of
      \S\ref{sec:mu2}. Guard (partial): the sweep \emph{retains} the
      previous iterate at any point that does not clear
      (\code{recursive\_main.\_sweep}, \code{keep\_tol}) --- which stops a
      single blow-up but cannot stop the contamination.
\item \textbf{Reaching a deep crisis from a calm start.} The post-default
      regime is far from the steady-state cold start, so a solver launched
      from the steady state may not find it. Guard: a \emph{homotopy} ---
      lower the recovery rate in steps $0.85\to0.70\to0.55\to$ target,
      re-solving each (\code{recursive\_experiment.solve\_recursive}).
\item \textbf{A slow outer loop.} The rules are found by \emph{time
      iteration} --- guess the rules, use them as next period's behaviour,
      re-solve, repeat. Its convergence speed is tied to the model's
      persistence ($\approx0.98$ for capital), so it crawls and can wobble.
      Guards: damping (blend new and old), and measuring convergence on rule
      \emph{values}.
\end{enumerate}


\section{The branch computation}
\label{sec:branch}

Now the piece you asked about: how ``the branch'' is computed. The word
refers to how the model handles the fork in the road each period --- next
period the sovereign either defaults ($d'=1$) or does not ($d'=0$) --- inside
the expectations that appear in the banks' optimality conditions. Your repo
contains \emph{two} ways of doing this, and the difference between them is
the whole reason the recursive solver exists.

\subsection{The economic object: an expectation over a fork}

A banker's optimality condition (an Euler equation) says, roughly, ``the
price of an asset today equals the expected discounted value of its payoff
tomorrow.'' The expectation runs over next period's shocks \emph{and} over
the default fork. For any payoff $X'$,
\begin{equation}
\E\big[X'\big]
=\underbrace{\sum_{\text{shock draws }j} w_j}_{\text{average over shocks}}
\Big[\underbrace{(1-p^d)\,X'^{(d'=0)}}_{\text{no-default branch}}
+\underbrace{p^d\,X'^{(d'=1)}}_{\text{default branch}}\Big],
\label{eq:expect}
\end{equation}
where the probability of default is a logistic function of the risk factor
(\code{state\_grid.default\_prob}),
\begin{equation}
p^d(s)=\frac{1}{1+e^{-s}}.
\label{eq:pd}
\end{equation}
At your calibration (\code{s\_process\_params}) the resting value is
$s^\star=-6.91$, giving $p^d\approx0.1\%$, and a two-standard-deviation risk
shock lifts $s$ to about $-3.9$, i.e. $p^d\approx2\%$ (Figure~\ref{fig:branch},
right). The crucial modelling point is that \eqref{eq:expect} is a genuine
\emph{probability-weighted average over two possible futures} --- a
distribution, not a single representative number. Raising $s$ shifts weight
onto the low-payoff default branch.

\subsection{The approximation that was removed}

An earlier perfect-foresight version priced this fork with a single
\emph{representative} ``feared default'' economy: one stand-in default economy
computed once and reused at every date, its low payoffs paired with a high
marginal valuation to price a risk premium. It was a reasonable-sounding
shortcut, but it collapsed the distribution to a single point, and in this
model that got the sign \emph{wrong}. With only one frozen default economy,
capital ended up looking like the safe place to hide, so banks levered
\emph{into} capital when risk rose and output \emph{expanded} --- the opposite
of the data and of Bocola's result. That approximation has been removed; the
two-branch quadrature below is the only way risk is priced.

\subsection{The right way (two-branch quadrature, \code{point\_map.py})}

The recursive solver evaluates \eqref{eq:expect} honestly. The ``average over
shocks'' is done by \emph{Gauss--Hermite quadrature} --- the standard, exact
tool for averaging against a bell curve. With the technology and spending
shocks held fixed for this experiment, the only continuous shock is the risk
innovation $\epsilon_s$, and $7$ Gauss--Hermite nodes suffice
(\code{expectations.gh\_nodes}). Each node is pushed one step forward through
the AR(1) law for $s$ to a next-period state; then the default fork doubles
it into a no-default and a default copy. The whole thing lives in
\code{point\_map.point\_residuals}: the helper \code{\_cont(d\_next)}
evaluates next period's rules in regime $d'$, and \code{\_E(v0,v1)} takes the
probability-weighted average (Figure~\ref{fig:branch}, left):
\begin{verbatim}
w_nd, w_def = wq * (1.0 - pd), wq * pd     # shock weights x branch probs
def _E(v0, v1):                            # two-branch expectation
    return float(np.dot(w_nd, v0) + np.dot(w_def, v1))
\end{verbatim}
The single most important idea is what \code{\_cont} does: the ``default
economy'' is \emph{not} a separate object bolted on --- it is \textbf{the same
decision rules, evaluated at a reachable next-period state}, in the $d'=1$
coefficient set. The bad state is somewhere the rules already describe; you
just arrive there with probability $p^d$. Three things differ between the two
branches, all visible in \code{point\_map.py}:

\begin{itemize}
\item \textbf{Which rules.} Each rule has a separate coefficient set per
      regime, \code{coef[name][d]}; the default branch reads the $d'=1$ set,
      the no-default branch the $d'=0$ set.
\item \textbf{The bond payoff.} In the default branch the $D$-bond is
      written down to its recovery value (\code{recovery\_rate\_D}); this
      haircut, weighted by $p^d$, is what pulls the bond price down when risk
      rises.
\item \textbf{The discount factor.} The no-default branch discounts with the
      bank's usual kernel $\Omega^0=\beta_{\text{inter}}(f+(1-f)\alpha')$;
      the default branch uses an income-based factor
      $\Omega^1=\Lambda^1(f+(1-f)\alpha')$ with
      $\Lambda^1=\beta_{\text{inter}}(Y^{d'=1}/Y^{d'=0})^{-\sigma}$ --- higher
      marginal value when default-state income is low. This is what makes the
      low default payoff \emph{extra} costly.
\end{itemize}

\begin{figure}[h]
\centering
\includegraphics[width=\textwidth]{branch_global.pdf}
\caption{\emph{Left:} the two-branch quadrature in \code{point\_map.py}. The
current state goes forward through the AR(1) for $s$ to $7$ Gauss--Hermite
draws; each splits into a no-default continuation (weight $1-p^d$, ordinary
discounting) and a default continuation (weight $p^d$, bond haircut and
income discounting); the expectation is their weighted sum. \emph{Right:}
the logistic default probability, with the resting point $s^\star$ ($0.1\%$)
and a $+2\sigma$ risk shock ($2\%$) marked.}
\label{fig:branch}
\end{figure}

\demohead{The two-branch expectation with numbers.}
Take a bond that pays $1$ next period if there is no default and, after a
$55\%$ haircut, $0.45$ if there is (recovery $0.45$). At the resting risk
level $p^d=0.001$ the expected payoff is
$(1-0.001)\cdot1+0.001\cdot0.45=0.99945$ --- essentially $1$. Now hit the
economy with the $+2\sigma$ shock, $p^d=0.02$:
\[
\E[\text{payoff}]=(1-0.02)\cdot1+0.02\cdot0.45=0.989.
\]
The \emph{average} payoff falls about $1.1\%$. But the price falls by
\emph{more}, because the default branch also carries the higher discount
factor $\Omega^1>\Omega^0$. Suppose $\Omega^0=1$ and, with income down in the
default state, $\Omega^1=1.3$. The bond-Euler numerator the code forms,
$\E[\Omega\cdot\text{payoff}]$, goes from
$0.98\cdot1\cdot1+0.02\cdot1.3\cdot0.45=0.9917$ at $p^d=0.02$ against a
no-risk value of $1$ --- and the same rise in $p^d$ lowers bank net worth,
which raises the multiplier $\mu$ in the denominator of the bond-price
formula \code{Q\_bD\_new = E\_Om\_payD / (E\_Om\_D*(1+rdep) + lambda*mu)}.
Both moves push $Q_{bD}$ down. \emph{That} is the pass-through: higher $s\Rightarrow$
more weight on the haircut branch $\Rightarrow$ lower bond price
$\Rightarrow$ mark-to-market loss on bank balance sheets $\Rightarrow$ tighter
constraint and wider lending spreads $\Rightarrow$ lower investment and
output. With the distribution represented honestly, the sign comes out right:
elevated default risk \emph{lowers} $Y_D$ and $C_D$, persistently
(the result read off by \code{recursive\_experiment.impact\_table} and
\code{persistence\_irf}).
\demofoot


\section{Why the solve fails at \texorpdfstring{$\mu=2$}{mu=2} but works at
\texorpdfstring{$\mu=1$}{mu=1}}
\label{sec:mu2}

This is the question you flagged. The recursive experiment currently runs on
the $\mu=1$ grid by default (\code{solve\_recursive(\dots, mu=1)}), and the
code comment records the reason bluntly: at $\mu=2$ ``the $d=1$ regime does
not converge on the wide box (joint residual $\sim0.7$, sign flips
expansionary).'' Here is what is going on and why it is not a bug you can
patch away --- it is the same global-basis obstacle documented for the
standalone Bocola solve in \code{docs/bocola2016\_replication.md}.

\subsection{The symptom}

Recall the two grids at six states: $\mu=1$ has $13$ points, $\mu=2$ has
$85$. When you solve both default regimes:

\begin{itemize}
\item \textbf{$\mu=1$ converges} to a small residual and gives the
      economically correct, contractionary sign. Its weakness is only that,
      being nearly separable (\S\ref{sec:smolyak}), it cannot represent how
      the responses interact across states --- so its magnitudes are
      ``indicative,'' but its \emph{sign and mechanism are right}.
\item \textbf{$\mu=2$ does not converge.} The no-default ($d=0$) block is
      fine, but the joint two-regime solve leaves a worst residual around
      $0.7$ --- six orders of magnitude above the $10^{-6}$ target --- and,
      worse, the implied risk response flips to \emph{expansionary}, the very
      artifact the recursive solver was built to eliminate.
\end{itemize}

\subsection{Why $\mu=1$ gets away with it}

Look again at which points each grid uses (\S\ref{sec:smolyak}). The $\mu=1$
grid is just the \emph{axes}: the centre, plus a step along each state one at
a time. Every one of those $13$ points sits close to the steady state (the
box bands are tight, $2\%$--$8\%$), so every point is economically
well-behaved --- positive net worth, a constraint that is at most gently
binding, a nonlinear solve that clears cleanly. There are no ``two things
extreme at once'' points. The homotopy walks the $d=1$ regime down to the
full haircut without ever meeting a state it cannot solve. Clean inputs
everywhere give a clean fit and the loop converges. The price, as
\S\ref{sec:connect-conv} said, is that a near-separable grid cannot represent
how the responses interact across states --- so the magnitudes are
indicative, though the sign and mechanism are right.

\subsection{The potential reasons $\mu=2$ will not converge}
\label{sec:mu2-reasons}

Going from $\mu=1$ to $\mu=2$ adds exactly the points $\mu=1$ omits: the
\emph{cross-corner} nodes, where two states are extreme at the same time (the
$(2,2)$ block of the worked example, the four corners of
Figure~\ref{fig:blocks}). In six dimensions these are combinations like
\emph{low capital together with high risk}, or \emph{high deposit obligation
together with low surviving debt}. Adding them sets off not one problem but a
stack of reinforcing ones. Here they are in plain terms, ordered roughly by
how much the code's own diagnostics blame them.

\begin{enumerate}
\item \textbf{The constraint barely binds, so there are two nearby answers
      and the fine grid grabs the wrong one.} \emph{(The deepest reason; the
      diagnostics confirm it.)} At the steady state the leverage multiplier
      is tiny, $\mu_{\text{ss}}\approx0.02$ --- the constraint is only just
      switched on, sitting right next to the $\max\{\cdot,0\}$ kink. That
      close to the switch, the model has \emph{two} valid equilibria almost
      on top of each other: a \emph{binding} one, where a bank that cannot
      lever up responds to risk by de-risking (output falls, the right sign),
      and a \emph{slack} one, where the bank substitutes into capital
      (output rises, the wrong sign). They are separated by a sliver. The
      coarse $\mu=1$ grid is too blunt to ``see'' the slack answer, so the
      solver stays on the binding one. The finer $\mu=2$ grid --- more points,
      packed closer, with a root-finder launched from more places --- keeps
      slipping onto the slack branch. That slip \emph{is} the expansionary
      sign it reports. \emph{Plainly:} two correct-looking answers sit almost
      on top of each other, and the fine grid keeps grabbing the wrong one.

\item \textbf{A few corners cannot be solved at all, and a global fit spreads
      their error everywhere.} \emph{(Confirmed: the code retains failed
      points, and the fit is global.)} In the already-defaulted regime the
      $55\%$ haircut has crushed net worth on arrival. At a cross-corner that
      \emph{also} has, say, high deposit obligations, the bank is near
      insolvency --- $\mu$ is pinned at its cap, net worth is tiny or
      negative --- and the pointwise solver finds no sensible equilibrium.
      The code sensibly \emph{retains} the previous value there rather than
      crash (\code{\_sweep}, \code{keep\_tol}). But the Chebyshev basis is
      \emph{global}: every coefficient is a blend of \textbf{all} $85$ point
      values (the single $\Phi c=f$ solve). So a handful of frozen, wrong
      corner values leak into every coefficient, and therefore into the
      fitted surface \emph{everywhere} --- including the calm interior where we
      read the answer. The interior then ``solves'' against a poisoned
      forecast, and the joint residual can never fall below $\sim0.7$.
      \emph{Plainly:} one wrong entry in a shared spreadsheet formula throws
      off every cell.

\item \textbf{The point solver is local and starts from one guess.}
      \emph{(Structural.)} Each grid-point solve is a local root-finder
      (\code{scipy.root}) launched from the steady-state / warm-start guess.
      The deep-recession $d=1$ solution can have a tiny ``catchment'' --- you
      only land on it if you start nearby. $\mu=2$ adds many corner states far
      from any good guess, so more of them find no root, or the wrong one.
      \emph{Plainly:} you are feeling for a light-switch in the dark from
      where you happen to stand, and the fine grid adds rooms you cannot reach
      from there.

\item \textbf{A polynomial cannot fold sharply, so it ripples at the crease.}
      \emph{(Structural; \S\ref{sec:accuracy}.)} The true rules have a kink
      where the constraint turns on. A smooth polynomial forced to bend
      sharply overshoots and wiggles nearby (the Gibbs ripples). Those wiggles
      can nudge a neighbouring point to the wrong side of the
      $\max\{\cdot,0\}$ switch, flipping it between the binding and slack
      regimes from one iteration to the next, so it never settles. $\mu=2$ has
      more points near the crease, so more of them wobble. \emph{Plainly:}
      drawing a sharp corner with a springy ruler --- press harder and it
      bounces.

\item \textbf{The two regimes are solved against each other, so errors echo.}
      \emph{(Structural.)} The $d=0$ and $d=1$ rules share the continuation:
      tomorrow's default-state rules feed today's no-default expectation and
      the other way round. A single bad $d=1$ corner therefore contaminates
      $d=0$ too, and back again. The joint solve is stiffer than either regime
      alone --- and it is precisely the joint step where the residual sticks.
      \emph{Plainly:} two people copying each other's homework; one mistake
      spreads both ways.

\item \textbf{The box cannot be both wide enough and clean.}
      \emph{(Structural.)} Capital is almost a random walk (persistence
      $\approx0.98$), so the states the model truly visits form a wide,
      tilted cloud. A box wide enough to hold it has economically impossible
      corners; a box narrow enough to avoid them clips the continuation.
      Whatever box you choose, $\mu=2$ populates its corners --- and in $d=1$
      those corners are the near-insolvent states of reason 2.
      \emph{Plainly:} the net must be big to catch the fish, but every extra
      metre of net drags in junk.

\item \textbf{Retaining failed points is a bandage, not a cure.}
      \emph{(Structural; follows from 2.)} Freezing a point that will not
      solve keeps the run alive, but a frozen value is still a value that the
      global fit --- and so the interior's forecast --- treats as truth. It
      stops a crash; it does not stop the contamination.

\item \textbf{The outer loop is slow and can limit-cycle at the crease.}
      \emph{(Structural.)} Time iteration's speed is tied to the model's
      persistence ($\approx0.98$), so it crawls; and at the kink the affected
      points can oscillate between binding and slack across iterations without
      ever settling. Damping softens this but does not remove it.
      \emph{Plainly:} a sluggish thermostat that keeps overshooting.
\end{enumerate}

\noindent Reasons 1 and 2 are the load-bearing ones --- the knife-edge
multiplicity and the global-fit corner poisoning --- with 3--8 as reinforcing
structural difficulties. And they are not independent: the poisoned corners
(2) push interior points toward the slack branch (1), and the crease wiggles
(4) do the same, so the failure feeds on itself.

\subsection{Why the sign flips, specifically}

The flip to an \emph{expansionary} response is not a separate bug --- it is
the fingerprint of reasons 1 and 2. Recall from \S\ref{sec:branch} that the
two-branch mechanism gives the contractionary sign \emph{only when the
default-branch continuation is a genuine recession}: the default state must
have low income and a low capital return, so that the high discount factor
$\Omega^1$ multiplies a low payoff and makes risk costly. Both a slip onto
the slack branch (reason 1) and a poisoned $d=1$ continuation (reason 2)
destroy that: capital reappears as a safe haven, banks lever \emph{into} it
when risk rises, and output expands --- exactly the expansionary artifact the
recursive solver was built to avoid. So seeing the expansionary
sign is how you \emph{know} the solve has drifted onto the wrong branch or
poisoned its corners.

\subsection{What would actually fix it}

The levers line up with the reasons, and several are one-line changes the
code already supports.

\begin{itemize}
\item \textbf{Bind a little harder in the calibration --- attacks reason 1,
      the root fix.} \emph{(A direction the diagnostics point to.)} The whole
      knife-edge exists because $\mu_{\text{ss}}\approx0.02$ is barely
      positive. Lowering $\lambda_K$ (or tightening the leverage target) so
      that $\mu_{\text{ss}}$ sits comfortably around $0.05$--$0.10$ pushes the
      binding equilibrium away from the slack one; the binding branch becomes
      the unique, robust answer and the fine grid can no longer slip onto the
      wrong one. The cost is a small step away from Bocola's very small
      estimated friction --- a modelling choice, not a bug.
\item \textbf{Read the rules instead of re-solving --- attacks reason 1 at
      read time.} The impact and IRF routines already evaluate the converged
      rules at their own policy values (\code{read\_at}) rather than launching
      a fresh root-finder, precisely so the read stays on the binding branch
      instead of slipping to the slack one. Keeping every downstream read on
      the rules (never re-solving) is the cheap safeguard.
\item \textbf{Refine only the states that move --- attacks reasons 2 and 6,
      cheaply.} \code{build\_state\_box} accepts \code{mu\_vec}: keep the
      near-fixed capital dimensions at level 1 and give level 2 only to
      $(s,P,B)$, the states where the constraint boundary actually moves. This
      buys the cross-terms you want \emph{without} creating the low-capital
      corners that blow up in $d=1$. It is the cheapest first experiment, and
      the hook is already there.
\item \textbf{Use a local basis --- attacks reason 2 at its root.} Replace the
      global Chebyshev basis with a local one (splines / finite elements), so
      a hard corner's error stays local instead of leaking through the fit.
      This is the route the Bocola replication notes single out, having hit
      the identical wall (there the residual floored around
      $5\times10^{-3}$ on a rotated grid; here, on the wider two-country
      $d=1$ box, it is far worse at $\sim0.7$).
\item \textbf{Shape the box to the cloud --- attacks reason 6.} A
      PCA-rotated box (used in the standalone Bocola solve) cuts the
      infeasible-corner fraction sharply, so fewer nodes sit on impossible
      states in the first place.
\item \textbf{Extend the homotopy --- attacks reason 3.} The recovery-rate
      ladder that makes the near-steady-state $d=1$ solve work could be paired
      with a \emph{box-widening} homotopy, growing $\mu=2$'s reach only as the
      solution stabilizes.
\end{itemize}

\noindent Honest summary: $\mu=1$ gives a converged, correctly-signed answer
whose magnitudes are indicative; a converged $\mu=2$ answer is not a matter of
more iterations but of \emph{moving off the knife-edge} (calibration) and/or
\emph{decoupling the hard corners from the interior} (a different
representation). The two first things to try are the anisotropic
\code{mu\_vec} grid and the bind-harder calibration, because both are
one-line changes the code already supports.


\section{Putting it together}
\label{sec:together}

The full solve (\code{recursive\_experiment.solve\_recursive}) chains the
pieces:
\begin{enumerate}
\item build the Smolyak box around the steady state and cold-start every rule
      at its steady-state value (\code{RuleSet.from\_ss});
\item calibrate the household budget anchors so the steady state is an exact
      rest point (\code{calibrate\_household\_anchors});
\item time-iterate the no-default ($d=0$) rules to convergence;
\item warm-start the default ($d=1$) rules from $d=0$ and walk the recovery
      rate down by homotopy;
\item refine both regimes jointly;
\item read the pass-through: the impact response as $p^d$ rises
      (\code{impact\_table}) and the decay of a one-off risk shock
      (\code{persistence\_irf}).
\end{enumerate}
Each sweep of the time iteration (\code{recursive\_main.time\_iteration})
freezes the current rules as next period's behaviour, solves the seven
market-clearing unknowns at every grid point and regime, reads off the
banker valuations and household aggregates, then damps and refits. It is the
per-point image of the same equations the perfect-foresight solver stacks ---
the difference is that here the answer is a reusable \emph{policy}, and the
expectation over the default fork is genuine.


\section*{Summary}
\addcontentsline{toc}{section}{Summary}

Three tools, one caveat. \emph{Chebyshev interpolation} turns each policy
function into a short, stable list of coefficients, near-optimal for smooth
targets but explosive if you extrapolate and slow at kinks.
\emph{Smolyak grids} carry that into six dimensions for a handful of points
instead of thousands, by keeping only the low-total-level combinations of
nested Chebyshev levels. The \emph{two-branch expectation} in
\code{point\_map.py} prices sovereign risk honestly, as a probability weight
on a default continuation that is the same rules read at a reachable state
--- which is what gets the sign right where the earlier approximation got it
wrong. The caveat runs through all of it: a global polynomial basis couples
every collocation point to every other, so the method is only as good as the
worst state you ask it to represent. That is why the box is kept narrow, why
every evaluation is clipped, why failed points are retained rather than
trusted --- and, ultimately, why $\mu=2$ stalls while $\mu=1$ converges: the
extra cross-corner points $\mu=2$ introduces are states the $d=1$ regime
cannot solve, and a global basis cannot stop their contamination from
spreading.


\section*{Where to read the code}
\addcontentsline{toc}{section}{Where to read the code}
\begin{center}
\small
\begin{tabular}{@{}ll@{}}
\toprule
Concept & File / function (\code{code/global/}) \\
\midrule
Chebyshev recurrence, nodes & \code{solver\_recursive/state\_grid.py}: \code{chebyshev\_basis\_1d}, \code{\_level\_points} \\
Smolyak combination rule & \code{solver\_recursive/state\_grid.py}: \code{\_multi\_indices}, \code{SmolyakGrid} \\
Fit / eval (cached LU solve) & \code{solver\_recursive/state\_grid.py}: \code{fit}, \code{eval}, \code{clip} \\
State box, anisotropy & \code{solver\_recursive/state\_grid.py}: \code{build\_state\_box}, \code{mu\_vec} \\
Per-regime rule storage & \code{solver\_recursive/decision\_rules.py}: \code{RuleSet}, \code{SOLVE7} \\
Two-branch expectation & \code{solver\_recursive/point\_map.py}: \code{\_cont}, \code{\_E} \\
Closed-form multiplier $\mu$ & \code{solver\_recursive/point\_map.py} (Bocola Prop.\ 1) \\
Default probability & \code{solver\_recursive/state\_grid.py}: \code{default\_prob}, \code{s\_process\_params} \\
Time iteration, retain guard & \code{solver\_recursive/recursive\_main.py}: \code{time\_iteration}, \code{\_sweep} \\
Homotopy, $\mu=1$ default, IRFs & \code{solver\_recursive/recursive\_experiment.py} \\
OMT/TPI three-branch, activation & \code{solver\_recursive/tpi\_recursive\_experiment.py} \\
\bottomrule
\end{tabular}
\end{center}


\section*{References and further reading}
\addcontentsline{toc}{section}{References and further reading}
\small
\begin{itemize}
\item L.\ Bocola (2016), ``The Pass-Through of Sovereign Risk,''
      \emph{Journal of Political Economy} 124(4) --- the economics the
      two-branch expectation implements.
\item D.\ Krueger and F.\ Kubler (2004), \emph{J.\ Economic Dynamics \&
      Control} 28 --- the nested Chebyshev--Smolyak construction the code
      follows.
\item K.\ Judd, L.\ Maliar, S.\ Maliar and R.\ Valero (2014), \emph{J.\
      Economic Dynamics \& Control} 44 --- anisotropic Smolyak levels (the
      \code{mu\_vec} idea).
      \url{https://bfi.uchicago.edu/wp-content/uploads/Judd-Maliar-Valero-1.pdf}
\item L.\ N.\ Trefethen (2013), \emph{Approximation Theory and Approximation
      Practice}, SIAM --- Chebyshev interpolation, the Runge phenomenon, and
      spectral convergence, at a readable level.
\item J.\ Brumm and S.\ Scheidegger (2017), \emph{Econometrica} 85 ---
      adaptive/local sparse grids, the corner-decoupling route out of the
      $\mu=2$ obstacle.
\item Internal: \code{docs/bocola2016\_replication.md} --- the identical
      global-basis corner obstacle, analyzed for the standalone solve.
\end{itemize}

\end{document}
